\section{Architecture du pipeline de bout en bout}
\label{sec:architecture-pipeline}

Cette partie décrit le chemin effectif d'un audit LANCE, depuis la sélection
d'un mode d'exécution jusqu'aux fichiers déposés dans le dossier du run, en
s'appuyant exclusivement sur le code de la branche \code{dev/1}. Les
affirmations renvoient au module et, lorsque c'est utile, à la fonction qui les
implémente ; les comportements que le code ne permet pas d'établir sont
signalés comme tels plutôt que comblés par inférence.

\subsection{Vue d'ensemble et sélection du mode}

L'application ne porte pas la logique d'audit : elle valide la demande, garantit
un seul run actif par instance, diffuse les événements SSE et relit les
artefacts (\code{docs/architecture/run-modes.md}). C'est \code{make\_runner}
(\code{src/agent/run\_modes.py}) qui choisit la stratégie d'exécution. Pour
\code{lance} et \code{rules}, il instancie un \code{Pipeline} et l'encapsule
dans \code{LanceRunner}, adaptateur qui réutilise l'objet pipeline et son API
publique. Pour \code{vanilla} et \code{scripted}, il construit un
\code{SimpleRunner}, classe autonome qui ne partage avec le pipeline que la
frontière d'exécution des outils, le journal de preuves, le suivi de coût et la
réservation du laboratoire.

\begin{table}[H]
\centering
\caption{Modes d'exécution et frontière correspondante.}
\label{tab:run-modes}
\small
\begin{tabular}{@{}llll@{}}
\toprule
Mode & Moteur & Modèle & Effets \\
\midrule
\code{lance} & \code{Pipeline} (6 phases) & requis & scénario, déploiement, évaluation \\
\code{rules} & \code{Pipeline} (phases 3--4) & aucun & scénario + audit déterministe \\
\code{vanilla} & \code{SimpleRunner} & requis & sonde Nmap bornée sur CIDR \\
\code{scripted} & \code{SimpleRunner} & aucun & procédures versionnées \\
\bottomrule
\end{tabular}
\end{table}

Le mode \code{rules} n'est pas un second moteur : \code{make\_runner} fixe
\code{decision\_policy="rules"}, \code{phases=[3, 4]} et
\code{execution\_profile="full"}, et refuse tout fournisseur, réseau direct ou
budget monétaire. Les modes \code{vanilla} et \code{scripted} rejettent au
contraire \code{scenario\_id} et tout réglage propre au harnais. Le qualificatif
\code{full}/\code{compact} désigne l'orchestration du pipeline, jamais le type
de runner.

\subsection{Les six phases et leurs livrables}

Le registre déclaratif \code{src/agent/registry.py} (\code{AGENTS}) attribue à
chaque phase un modèle de prompt, un groupe d'outils, un validateur et un nom de
livrable. Le routage est explicite dans \code{src/agent/phases/registry.py} :
\code{PHASE\_MODULES} associe les numéros 1 à 6 à des modules fixes, sans import
dynamique.

\begin{table}[H]
\centering
\caption{Phases, modules, livrables principaux et artefacts associés.}
\label{tab:phases}
\small
\begin{tabular}{@{}lp{2.7cm}p{3.2cm}p{4.05cm}@{}}
\toprule
Phase & Module & Livrable principal & Artefacts associés \\
\midrule
1 & \code{graph/run.py} & \code{01\_graph\_analysis.md} & \code{01\_graph\_evidence.json} \\
2 & \code{recon/run.py} & \code{02\_recon.md} & \code{02\_recon\_evidence.json}, \code{02\_topology\_edges.json} \\
3 & \code{analysis/run.py} & \code{03\_vuln\_analysis.json} & \code{03\_vuln\_analysis\_raw.json}, \code{03\_phase3\_status.json}, \code{03\_device\_*} \\
4 & \code{verification/run.py} & \code{04\_exploitation.json} & diagnostics \code{execution\_errors} \\
5 & \code{intrusion/run.py} & \code{05\_intrusion.json} & \code{05\_intrusion\_context.json} \\
6 & \code{report/run.py} & \code{06\_report.md} & \code{06\_report\_analysis.md}, fiches par section \\
\bottomrule
\end{tabular}
\end{table}

La phase 1 projette les résultats des outils de graphe en un journal de faits
indépendant du discours du modèle
(\code{GraphPhase.\_build\_graph\_evidence\_projection}), et le runner refuse
une soumission dont la couverture par appareil est incomplète
(\code{src/agent/core/runner.py}). La phase 2 laisse le modèle choisir les
sondes, mais \code{02\_recon.md} est rendu depuis le journal d'outils, non
recopié par le modèle. La phase 3 s'exécute en trois temps :
\code{scan.scan\_phase} (observations déterministes),
\code{devices.analyze\_devices} (analystes par appareil) puis
\code{run.aggregate\_phase}, qui normalise, déduplique et projette les constats
et écrit \code{03\_vuln\_analysis.json}. La phase 4 lance des micro-agents par
vulnérabilité et agrège leurs résultats. La phase 5 produit son bilan depuis le
journal : le modèle ne fournit qu'un marqueur de fin, et le statut reste
\code{objectives\_status=not\_certified}. La phase 6 est particulière : son
\code{AgentConfig} ne déclare aucun outil et \code{max\_turns=1} ; le rapport
est assemblé de façon déterministe à partir de fiches, le modèle ne rédigeant
qu'une note d'analyste bornée (\code{06\_report\_analysis.md}).

\subsection{Graphe et contexte de run}

Le graphe est construit en deux couches. \code{src/loader.py} lit le YAML
d'infrastructure en dataclasses (\code{load\_yaml}) puis remplit un backend
(\code{build\_graph}, backend par défaut \code{NetworkXBackend}).
\code{src/graph\_backend.py} définit l'interface abstraite \code{GraphBackend}
--- \code{get\_device}, \code{get\_neighbors}, \code{find\_all\_paths},
\code{get\_attack\_surface}, \code{get\_graph\_stats}, \code{to\_dict} --- et son
implémentation NetworkX sur un \code{DiGraph}. Deux détails méritent d'être
soulignés : \code{get\_neighbors} renvoie l'union triée des successeurs et
prédécesseurs, et \code{find\_all\_paths} travaille sur une vue non orientée
(\code{to\_undirected}) afin de trouver des chemins dans les deux sens, y
compris à travers le backend abstrait.

L'état appartient au run. Chaque \code{Pipeline} crée un \code{RunContext}
immuable (dossier du run et \code{GraphContext}) ; \code{GraphContext} porte le
backend, l'infrastructure, les rapports CVE, les scores de risque, la topologie
de scénario ou de découverte, le graphe pondéré et le dernier delta. Les outils
de graphe retrouvent ce contexte par un \code{ContextVar}
(\code{src/agent/core/graph\_context.py}), sans modifier les arguments exposés
au modèle ; \code{graph\_bound} et \code{bind\_graph\_callable} capturent
explicitement le propriétaire, y compris pour les workers en threads, plutôt
que de compter sur la propagation implicite du contexte.

\subsection{La frontière de gouvernance des outils}

Tous les profils et tous les modes passent par \code{wrap\_tool}
(\code{src/agent/core/executor.py}). À chaque appel, la fonction enveloppée :
(1) vérifie que la preuve précédente a pu être archivée, l'arrêt coopératif et
les budgets via \code{check\_execution\_limits} ; (2) incrémente atomiquement un
compteur et refuse l'appel si \code{max\_tool\_calls} est atteint ; (3) attribue
une référence \code{tc-<uuid>} ; (4) applique les refus de portée ; (5) exécute
l'outil ; (6) archive systématiquement une entrée dans \code{tool\_calls.jsonl},
y compris pour un refus ou une exception.

\begin{table}[H]
\centering
\caption{Contrôles appliqués par \code{check\_execution\_limits} et \code{wrap\_tool}.}
\label{tab:executor}
\small
\begin{tabular}{@{}p{2.6cm}p{4.6cm}p{6.0cm}@{}}
\toprule
Contrôle & Mécanisme & Conséquence d'un échec \\
\midrule
Intégrité des preuves & drapeau \code{\_evidence\_integrity\_failed} & \code{EvidenceWriteError}, run non vérifiable \\
Arrêt utilisateur & \code{\_stop\_event} (portée expérimentale) & \code{RunStopped} \\
Budget de temps & \code{max\_duration\_s} via \code{time.monotonic()} & \code{BudgetExceeded} \\
Budget monétaire & \code{tracker.check\_budget()} & \code{BudgetExceeded} \\
Budget d'appels & compteur sous \code{\_artifact\_log\_lock} & \code{BudgetExceeded} \\
Périmètre & \code{\_intrusion\_scope\_violation} & refus JSON archivé \\
Exécution libre & \code{python\_exec} & refus \code{unverifiable\_execution\_scope} \\
\bottomrule
\end{tabular}
\end{table}

Deux propriétés sont importantes pour l'interprétation des résultats. D'une
part, un échec d'écriture du journal n'est jamais ignoré : il est converti en
\code{EvidenceWriteError}, ce qui interrompt le run plutôt que de laisser des
actions non tracées. D'autre part, la limite de coût de \code{check\_budget}
porte sur l'usage observé (\code{total\_cost()}), pas sur une réservation des
appels en vol ; le docstring du code le dit explicitement. Le périmètre est
contrôlé sur les champs de destination (\code{ip}, \code{host},
\code{target}, \code{url}...) contre le \code{target\_subnet} du run ; une
action sans destination vérifiable est refusée en amont (\code{python\_exec}),
et l'en-tête du module précise qu'il s'agit de gardes applicatives, non d'un
bac à sable système.

\subsection{Artefacts et suivi de coût}

Le dossier du run est créé de façon exclusive sous \code{output/agent/}
(\code{Pipeline.\_\_init\_\_}) ; en cas de collision, un suffixe hexadécimal est
ajouté et aucun dossier antérieur n'est écrasé. \code{run\_meta.json} consigne
phase, profil, politiques, commit et limites ; \code{cost\_summary.json} est
écrit dans le bloc \code{finally} du pipeline, donc même sur exception. Le
contrat de collecte (\code{src/benchmark/artifacts.py}) définit une liste
positive de fichiers autorisés (\code{\_ALLOWED\_ROOT\_FILES}) et des motifs
\code{03\_device\_*.json} ; \code{create\_submission\_bundle} produit un zip muni
d'un manifeste SHA-256 et revérifie le manifeste après écriture. Côté agent, une
tentative de livrable est archivée dans \code{.attempts/} et consignée dans
\code{deliverable\_attempts.jsonl} ; la route
\code{GET /api/runs/\{id\}/download/zip} exclut les fichiers privés et les
entrées cachées.

Le suivi de coût (\code{src/agent/cost\_tracker.py}) résout le prix une fois par
phase (\code{\_resolve\_pricing}) et l'inscrit dans \code{PhaseUsage}, de sorte
qu'un run terminé ne peut pas être réévalué a posteriori. Chaque phase porte la
source du prix (\code{dynamic\_catalog}, \code{static\_catalog},
\code{subscription}, \code{default\_estimate}) et un indicateur
\code{cost\_is\_estimate}. \code{summary()} distingue explicitement la durée
cumulée des agents (\code{total\_agent\_duration\_s}) du temps mural
(\code{wall\_clock\_duration\_s}) ; la première ne mesure pas le temps
réellement écoulé.

\subsection{Ce que l'architecture ne garantit pas}

L'architecture est isolée par run pour les livrables et le graphe, mais elle ne
l'est pas totalement : \code{src/agent/phases/README.md} indique que les
services de phase peuvent encore appeler le moteur commun et que la simultanéité
de scénarios sur un laboratoire partagé n'est pas garantie. \code{reserve\_lab}
n'est actif que si \code{LANCE\_LAB\_LOCK} est défini et n'est réentrant que dans
le même thread. Par ailleurs, sans paramètre explicite, le pipeline interactif
n'impose ni plafond d'appels ni délai d'analyse par équipement : les budgets du
tableau~\ref{tab:executor} sont donc conditionnels.

Trois points n'ont pas pu être établis par lecture du code. \emph{Un} : aucun
mécanisme d'archivage ou de compactage des dossiers de run n'a été trouvé dans
\code{src/} (aucune occurrence de \code{tarfile},
\code{shutil.make\_archive} ou équivalent) ; la conservation semble être celle du
système de fichiers, ce que la documentation ne décrit pas. \emph{Deux} :
\code{02\_recon\_evidence.json} est lu par \code{src/agent/scanner.py} et
\code{src/agent/core/runner.py}, mais son producteur principal n'a pas été
identifié ; l'affirmer serait une supposition. \emph{Trois} :
\code{provider\_events.jsonl} est consigné par le callback du fournisseur
(\code{core/runner.py}) via \code{provider\_diagnostics}, et le contenu exact des
champs n'a pas été vérifié ligne à ligne. Enfin, un statut \code{completed}
signifie que le runner a terminé son travail, jamais que la cible est couverte
exhaustivement ni qu'un accès a été prouvé
(\code{docs/architecture/run-modes.md}).
